\subsection{Definition of homotopies}

DEFINITION 1. --- Let G be a groupoid, let H be a quiver, and let \(\varphi\) and \(\varphi'\) be morphisms of quivers from H to G. A homotopy connecting \(\varphi\) to \(\varphi'\) is a map h from the vertex set of H to the set of arrows of G satisfying the following properties:

\begin{enumerate}
\def\labelenumi{(\roman{enumi})}
\item
  For every vertex a of H, the arrow \(h(a)\) has as its origin \(\varphi(a)\) and as its target \(\varphi'(a)\) ;
\item
  For every arrow \(f\) of H, with origin \(a\) and target \(b\) , we have \(\varphi(f)h(b)=h(a)\varphi'(f)\) .
\end{enumerate}

We say that \(\varphi\) and \(\varphi'\) are homotopic if there exists a homotopy connecting \(\varphi\) to \(\varphi'\) .

Let G be a groupoid, let H be a quiver and let \(\varphi\) , \(\varphi'\) , \(\varphi''\) be morphisms of quivers from H to G. The map \(a \mapsto e_{\varphi(a)}\) is a homotopy connecting \(\varphi\) to \(\varphi\) . If \(h\) is a homotopy connecting \(\varphi\) to \(\varphi'\) , the map \(a \mapsto h(a)^{-1}\) is a homotopy connecting \(\varphi'\) to \(\varphi\) . Let \(h\) and \(h'\) be homotopies connecting \(\varphi\) to \(\varphi'\) and \(\varphi'\) to \(\varphi''\) respectively. For every vertex \(a\) of H, the arrows \(h(a)\) and \(h'(a)\) are composable. The map \(a \mapsto h(a)h'(a)\) is a homotopy connecting \(\varphi\) to \(\varphi''\) .

It follows that the relation " \(\varphi\) is homotopic to \(\varphi'\) " is an equivalence relation on the set of quiver morphisms from H to G.

Let G be a groupoid, let H be a quiver, and let \(\varphi\) and \(\varphi'\) be morphisms of quivers from H to G that are homotopic. By condition (i) of Definition 1, for every vertex \(a\) of H, the vertices \(\varphi(a)\) and \(\varphi'(a)\) belong to the same orbit of G.

Suppose moreover that H is a groupoid and let h be a homotopy connecting \(\varphi\) to \(\varphi'\) . The maps \(\operatorname{Orb}(\varphi)\) and \(\operatorname{Orb}(\varphi')\) deduced from \(\varphi\) and \(\varphi'\) by passage to orbits are therefore equal. For every vertex a of H and every arrow \(f \in H_{a}\) , we have \(\varphi(f) = h(a)\varphi'(f)h(a)^{-1} = \operatorname{Int}(h(a))(\varphi'(f))\) , by condition (ii) of Definition 1. In other words, the homomorphism \(\varphi_{a}\) is equal to \(\operatorname{Int}(h(a)) \circ \varphi_{a}'\) . In particular, if the homomorphism \(\varphi_{a}\) is injective (resp. bijective, resp. surjective), then so is the homomorphism \(\varphi_{a}'\) .

Remarks. --- 1) Let G, G' be groupoids, let H, H' be quivers, let \(u\colon\mathrm{H}'\to\mathrm{H}\) be a morphism of quivers and \(v\colon\mathrm{G}\to\mathrm{G}'\) be a morphism of groupoids. If morphisms of quivers \(\varphi,\varphi'\) from H to G are homotopic, then the morphisms of quivers \(v\circ\varphi\circ u\) and \(v\circ\varphi'\circ u\) from H' to G' are homotopic. More precisely, if \(h\) is a homotopy connecting \(\varphi\) to \(\varphi'\) , the map \(\mathrm{Fl}(v)\circ h\circ\mathrm{Som}(u)\) is a homotopy connecting \(v\circ\varphi\circ u\) to \(v\circ\varphi'\circ u\) .

\begin{enumerate}
\def\labelenumi{\arabic{enumi})}
\setcounter{enumi}{1}
\tightlist
\item
  Let G be a groupoid, let H be a quiver and let \(\varphi, \psi\) be morphisms of quivers from H to G. Denote by \(j\) the canonical morphism from H to Grp(H), and let \(\overline{\varphi}\) and \(\overline{\psi}\) be the morphisms of groupoids from Grp(H) to G such that \(\overline{\varphi} \circ j = \varphi\) and \(\overline{\psi} \circ j = \psi\) .
\end{enumerate}

Recall that \(\operatorname{Som}(\mathrm{H}) = \operatorname{Som}(\operatorname{Grp}(\mathrm{H}))\) .

A homotopy \(h\colon \operatorname{Som}(\mathrm{H})\to \operatorname{Fl}(\mathrm{G})\) connecting \(\varphi\) to \(\psi\) is a homotopy connecting \(\overline{\varphi}\) to \(\overline{\psi}\) .

